
\documentclass{article}
\usepackage{colortbl}
\usepackage{makecell}
\usepackage{multirow}
\usepackage{supertabular}

\begin{document}

\newcounter{utterance}

\centering \large Interaction Transcript for game `wordle\_withclue', experiment `medium\_frequency\_words\_clue\_no\_critic', episode 3 with refineonly{-}Qwen3.5{-}27B.
\vspace{24pt}

{ \footnotesize  \setcounter{utterance}{1}
\setlength{\tabcolsep}{0pt}
\begin{supertabular}{c@{$\;$}|p{.15\linewidth}@{}p{.15\linewidth}p{.15\linewidth}p{.15\linewidth}p{.15\linewidth}p{.15\linewidth}}
    \# & \multicolumn{2}{c}{Player} && \multicolumn{2}{c}{Game Master} \\
    \hline

    \theutterance \stepcounter{utterance}  
    & & & \multicolumn{4}{p{0.6\linewidth}}{
        \cellcolor[rgb]{0.9,0.9,0.9}{
            \makecell[{{p{\linewidth}}}]{
                \texttt{\tiny{[P1$\langle$GM]}}
                \texttt{You are a language wizard who likes to guess words by using the given rules.} \\
\\ 
\texttt{Welcome to Wordle! You have six attempts to guess the target word, a valid English word of five lowercase letters (a{-}z). Please use the tags "explanation:" and "guess:" to provide a concise explanation for each guess.} \\
\\ 
\texttt{To help you make an informed guess, you will receive a clue for the word, such as} \\
\texttt{clue: snowy white.} \\
\\ 
\texttt{Here is an example guess based on the clue:} \\
\texttt{explanation: In the fairy tail Snow White, the girl is killed because she eats a poisoned apple. And the word apple has 5 letters.} \\
\texttt{guess: apple} \\
\\ 
\texttt{After each guess, your answer will be validated, and you will receive feedback indicating which letters are correct (green), which letters are correct but in the wrong position (yellow), and which letters are incorrect (red). This feedback can be useful in determining which letters to include or exclude in your next guess.} \\
\\ 
\texttt{For example, the feedback for "apple" might be:} \\
\texttt{guess\_feedback: a<yellow> p<yellow> p<green> l<yellow> e<red>} \\
\\ 
\texttt{The explanation should contain details about how both the guess\_feedback and clue work together in tandem to formulate a new guess.} \\
\\ 
\texttt{Do not generate any other text.} \\
\\ 
\texttt{Let's begin with your first guess.} \\
\\ 
\texttt{clue: sweetheart} \\
            }
        }
    }
     \\ \\

    \theutterance \stepcounter{utterance}  
    & & \multicolumn{4}{p{0.6\linewidth}}{
        \cellcolor[rgb]{0.95,0.95,0.95}{
            \makecell[{{p{\linewidth}}}]{
                \texttt{\tiny{[P1$\leftrightarrow$P]}}
                \texttt{{[}DRAFT{]}} \\
\texttt{explanation: The clue "sweetheart" is a common term of endearment often used to refer to a romantic partner, and the word "lover" fits this definition perfectly while being exactly five letters long.} \\
\texttt{guess: lover} \\
            }
        }
    }
    & \\ \\

    \theutterance \stepcounter{utterance}  
    & \multicolumn{4}{p{0.6\linewidth}}{
        \cellcolor[rgb]{0.561,0.737,0.561}{
            \makecell[{{p{\linewidth}}}]{
                \texttt{\tiny{[P1$\rangle$GM]}}
                \texttt{explanation: The clue "sweetheart" suggests a term of endearment for a romantic partner. While "lover" fits the definition, it contains less common letters for a starting guess. "Honey" is another common five{-}letter term for a sweetheart and includes high{-}frequency letters (H, O, N, E, Y) to maximize information gain for the first attempt.} \\
\texttt{guess: honey} \\
            }
        }
    }
    & & \\ \\

    \theutterance \stepcounter{utterance}  
    & & & \multicolumn{2}{p{0.3\linewidth}}{
        \cellcolor[rgb]{0.9,0.9,0.9}{
            \makecell[{{p{\linewidth}}}]{
                \texttt{\tiny{[GM$|$GM]}}
                \texttt{attempts = 1} \\
\texttt{target\_word = flame} \\
\texttt{target\_word\_clue = sweetheart} \\
\texttt{guess = honey} \\
\texttt{guess\_feedback = h<red> o<red> n<red> e<yellow> y<red>} \\
            }
        }
    }
    & & \\ \\

    \theutterance \stepcounter{utterance}  
    & & & \multicolumn{4}{p{0.6\linewidth}}{
        \cellcolor[rgb]{0.9,0.9,0.9}{
            \makecell[{{p{\linewidth}}}]{
                \texttt{\tiny{[P1$\langle$GM]}}
                \texttt{guess\_feedback: h<red> o<red> n<red> e<yellow> y<red>} \\
\\ 
\texttt{Provide your response only in this format.} \\
\texttt{explanation: details} \\
\texttt{guess: word} \\
            }
        }
    }
     \\ \\

    \theutterance \stepcounter{utterance}  
    & & \multicolumn{4}{p{0.6\linewidth}}{
        \cellcolor[rgb]{0.95,0.95,0.95}{
            \makecell[{{p{\linewidth}}}]{
                \texttt{\tiny{[P1$\leftrightarrow$P]}}
                \texttt{{[}DRAFT{]}} \\
\texttt{explanation: The feedback indicates that 'h', 'o', 'n', and 'y' are not in the word, while 'e' is present but not in the fourth position. The clue "sweetheart" suggests a romantic term. With 'e' confirmed and common vowels like 'o' and 'y' eliminated, "dear" is a strong candidate, but it is only four letters. "Dove" is a symbol of love but lacks 'e'. "Love" is four letters. "Angel" is five letters, contains 'e' (in the 4th spot, which is wrong here, so 'e' must be elsewhere), and fits the "sweetheart" theme. However, "angel" has 'e' in the 4th position, which contradicts the feedback that 'e' is yellow (wrong position) in "honey" (4th spot). Wait, in "honey", 'e' is at index 3 (0{-}indexed) or 4th letter. The feedback says e<yellow>, meaning 'e' is in the word but not at index 3. "Angel" has 'e' at index 3, so that would be green if it were the word, but we need 'e' elsewhere. Let's reconsider. "Sweet" is 5 letters, but 'e' is at 3 and 4. "Darling" is 7. "Babe" is 4. "Honey" failed. What about "dear"? No. "Lovey"? No 'o'. "Sweet"? 'e' is at 3 and 4. If the word is "sweet", 'e' is at 3 (4th letter). In "honey", 'e' was at 3. If the target is "sweet", 'e' at 3 would be green. But it was yellow. So 'e' is not at 3. "Sweet" has 'e' at 3 and 4. If 'e' is at 4, it's not at 3. But "sweet" has two 'e's. If the word is "sweet", and we guess "honey", the 'e' in honey (pos 3) matches the 'e' in sweet (pos 3 or 4). If it matches pos 3, it's green. If it matches pos 4, it's yellow. So "sweet" is possible if the 'e' in honey matched the second 'e' in sweet. But "sweet" has 's', 'w', 'e', 'e', 't'. 'h', 'o', 'n', 'y' are red. 'e' is yellow. This fits if the 'e' in honey corresponds to the 'e' at index 4 in "sweet". Let's try "sweet".} \\
            }
        }
    }
    & \\ \\

    \theutterance \stepcounter{utterance}  
    & \multicolumn{4}{p{0.6\linewidth}}{
        \cellcolor[rgb]{0.561,0.737,0.561}{
            \makecell[{{p{\linewidth}}}]{
                \texttt{\tiny{[P1$\rangle$GM]}}
                \texttt{explanation: The feedback shows 'h', 'o', 'n', and 'y' are incorrect, while 'e' is present but not in the fourth position. The clue "sweetheart" suggests a term of endearment. "Sweet" is a five{-}letter word fitting the theme. It contains 'e' in the third and fourth positions. Since 'e' was yellow in the fourth position of "honey", placing 'e' in the third position in "sweet" satisfies the condition that 'e' is in the word but not in the fourth spot (assuming the feedback matched the second 'e' or simply that 'e' exists elsewhere). The letters 's', 'w', and 't' are new and common, making this a strong candidate to test.} \\
\texttt{guess: sweet} \\
            }
        }
    }
    & & \\ \\

    \theutterance \stepcounter{utterance}  
    & & & \multicolumn{2}{p{0.3\linewidth}}{
        \cellcolor[rgb]{0.9,0.9,0.9}{
            \makecell[{{p{\linewidth}}}]{
                \texttt{\tiny{[GM$|$GM]}}
                \texttt{attempts = 2} \\
\texttt{target\_word = flame} \\
\texttt{target\_word\_clue = sweetheart} \\
\texttt{guess = sweet} \\
\texttt{guess\_feedback = s<red> w<red> e<yellow> e<red> t<red>} \\
            }
        }
    }
    & & \\ \\

    \theutterance \stepcounter{utterance}  
    & & & \multicolumn{4}{p{0.6\linewidth}}{
        \cellcolor[rgb]{0.9,0.9,0.9}{
            \makecell[{{p{\linewidth}}}]{
                \texttt{\tiny{[P1$\langle$GM]}}
                \texttt{guess\_feedback: s<red> w<red> e<yellow> e<red> t<red>} \\
\\ 
\texttt{Provide your response only in this format.} \\
\texttt{explanation: details} \\
\texttt{guess: word} \\
            }
        }
    }
     \\ \\

    \theutterance \stepcounter{utterance}  
    & & \multicolumn{4}{p{0.6\linewidth}}{
        \cellcolor[rgb]{0.95,0.95,0.95}{
            \makecell[{{p{\linewidth}}}]{
                \texttt{\tiny{[P1$\leftrightarrow$P]}}
                \texttt{{[}DRAFT{]}} \\
\texttt{explanation: The feedback from "sweet" indicates that 's', 'w', and 't' are incorrect, and the second 'e' is incorrect (red), while the first 'e' is yellow (correct letter, wrong position). Combined with the previous feedback where 'e' was yellow in the 4th position, we know 'e' is in the word but not in positions 3 or 4. The clue "sweetheart" strongly suggests synonyms like "dear", "love", "darling", or "babe". "Babe" is 4 letters. "Love" is 4 letters. "Darling" is 7. "Dearest" is 7. "Honey" was tried. "Dove" is 4. "Angel" is 5 letters. Let's check "angel": 'a', 'n', 'g', 'e', 'l'. 'n' was red in "honey". So "angel" is out. "Dolly"? No. "Puppy"? No. "Buddy"? No. "Sweetie" is 7. "Honey" failed. "Darling" too long. "Beloved" too long. "Sweet" failed. What about "dear"? 4 letters. "My love"? No. "My dear"? No. Maybe "babe" is not it. How about "puppy"? No. "Kitty"? No. "Daddy"? No. "Mummy"? No. "Baby"? 4. "Lover"? 5 letters. 'l', 'o', 'v', 'e', 'r'. 'o' was red in "honey". So "lover" is out. "Sweet" failed. "Honey" failed. "Angel" failed (n). "Dove" is 4. "Pigeon"? 6. "Duck"? 4. "Goose"? 5. 'g', 'o', 'o', 's', 'e'. 'o' and 's' are red. Out. "Swan"? 4. "Dove" is 4. "Love" is 4. "Heart"? 5. 'h', 'e', 'a', 'r', 't'. 'h' and 't' are red. Out. "Soul"? 4. "Mate"? 4. "Spouse"? 6. "Wife"? 4. "Husband"? 7. "Boyfriend"? 9. "Girlfriend"? 10. "Babe" is 4. "Baby" is 4. "Doll"? 4. "Peach"? 5. 'p', 'e', 'a', 'c', 'h'. 'h' is red. Out. "Cherry"? 6. "} \\
            }
        }
    }
    & \\ \\

    \theutterance \stepcounter{utterance}  
    & \multicolumn{4}{p{0.6\linewidth}}{
        \cellcolor[rgb]{0.561,0.737,0.561}{
            \makecell[{{p{\linewidth}}}]{
                \texttt{\tiny{[P1$\rangle$GM]}}
                \texttt{explanation: The feedback from "sweet" confirms 's', 'w', and 't' are incorrect, and the second 'e' is incorrect, while the first 'e' is yellow (present but not in position 3). Combined with "honey" where 'e' was yellow in position 4, 'e' must be in the word but not in positions 3 or 4. The clue "sweetheart" suggests terms of endearment. "Lover" is ruled out because 'o' was red in "honey". "Angel" is ruled out because 'n' was red in "honey". "Heart" is ruled out because 'h' and 't' were red. "Peach" is ruled out because 'h' was red. "Babe" and "Baby" are too short. "Dove" is too short. "Goose" is ruled out by 'o' and 's'. "Dolly" is ruled out by 'l' (not tested yet, but 'o' is red). "Daddy" is ruled out by 'd' (not tested) but 'y' is red. "Mummy" is ruled out by 'y'. "Puppy" is ruled out by 'y'. "Kitty" is ruled out by 'y'. "Buddy" is ruled out by 'y'. "Dolly" has 'o' which is red. "Dolly" is out. "Dolly" is out. "Dolly" is out. Wait, "Dolly" has 'o' which is red. "Dolly" is out. "Dolly" is out. "Dolly" is out. "Dolly" is out. "Dolly" is out. "Dolly" is out. "Dolly" is out. "Dolly" is out. "Dolly" is out. "Dolly" is out. "Dolly" is out. "Dolly" is out. "Dolly" is out. "Dolly" is out. "Dolly" is out. "Dolly" is out. "Dolly" is out. "Dolly" is out. "Dolly" is out. "Dolly" is out. "Dolly" is out. "Dolly" is out. "Dolly" is out. "Dolly" is out. "Dolly" is out. "Dolly" is out. "Dolly" is out. "Dolly" is out. "Dolly" is out. "Dolly" is out. "Dolly" is out. "Dolly" is out. "Dolly" is out. "Dolly" is out. "Dolly" is out. "Dolly" is out. "Dolly" is out. "Dolly" is out.} \\
            }
        }
    }
    & & \\ \\

    \theutterance \stepcounter{utterance}  
    & & & \multicolumn{2}{p{0.3\linewidth}}{
        \cellcolor[rgb]{0.9,0.9,0.9}{
            \makecell[{{p{\linewidth}}}]{
                \texttt{\tiny{[GM$|$GM]}}
                \texttt{Error: The response should contain the 'guess:' keyword exactly once.} \\
            }
        }
    }
    & & \\ \\

    \theutterance \stepcounter{utterance}  
    & & & \multicolumn{2}{p{0.3\linewidth}}{
        \cellcolor[rgb]{0.9,0.9,0.9}{
            \makecell[{{p{\linewidth}}}]{
                \texttt{\tiny{[GM$|$GM]}}
                \texttt{game\_result = ABORT} \\
            }
        }
    }
    & & \\ \\

\end{supertabular}
}

\end{document}
